\chapter{Create the Store Jail}

The catalog API is now reachable on its private address. We can introduce the
caller: a user-facing Rails application named \texttt{store}. It receives no
inbound host port because the future \texttt{web} jail will be the system's
only external entry point.

\section{The proposed service}

Use the following values:

\begin{description}
  \item[starter] Blank;
  \item[jail name] \texttt{store};
  \item[FreeBSD release] \texttt{15.1-RELEASE};
  \item[private address] \texttt{10.17.89.25/24};
  \item[outbound policy] Internet only;
  \item[starts at boot] disabled; and
  \item[inbound exposure] none.
\end{description}

The address is the next unused recommendation after \texttt{database} and
\texttt{catalog}. Internet access is temporary support for package and gem
installation, not permission for inbound traffic.

\section{Exercise 9A: Create the blank jail}

Open BSDRunner Jails and choose \textbf{Create a Service Jail}, then select the
\textbf{Blank} starter. Review the familiar five concepts:

\begin{enumerate}
  \item reuse the existing Blackship project and ZFS storage;
  \item select \texttt{15.1-RELEASE};
  \item use \texttt{10.17.89.25/24};
  \item name the jail \texttt{store}, select \textbf{Internet only}, and leave
    \textbf{Start at boot} disabled; and
  \item verify that the plan contains no inbound port forward.
\end{enumerate}

Type \texttt{store} into the confirmation field and create the jail. Once the
operation completes, start it and open its console.

\section{Checkpoint: The blank jail}

Stop after confirming that the store jail is running at
\texttt{10.17.89.25}. Do not install Rails yet. The next exercise will first
probe the catalog service from this otherwise blank jail, proving the private
application-to-application route before adding another Ruby stack.

\section{Exercise 9B: Call the catalog from the store jail}

The catalog server must be running while we perform the probe. In the catalog
jail console, become the application user and start Puma on the jail's private
address:

\begin{lstlisting}[language=sh]
su -l catalog
cd /usr/local/www/catalog
bin/rails server --binding 10.17.89.24 --port 3000
\end{lstlisting}

Leave that console open. In the store jail console, first confirm that the TCP
port is reachable:

\begin{lstlisting}[language=sh]
nc -vz -w 3 10.17.89.24 3000
\end{lstlisting}

Then make an HTTP request with \texttt{fetch}, which is part of the FreeBSD
base system:

\begin{lstlisting}[language=sh]
fetch -qo - http://10.17.89.24:3000/products
\end{lstlisting}

The first command should report a successful TCP connection. The second should
print the catalog's JSON product array. Together they prove that the store jail
can route to the catalog jail and that the Rails API answers over HTTP.

\section{Checkpoint: The private API route}

Keep the store console open, but stop here before installing its Rails
application. The next exercise will build the user-facing application only
after its dependency has been observed from inside the correct caller jail.
